\documentclass[hidelinks,a4paper,12pt]{article}
\addtolength{\oddsidemargin}{-1.cm}
\addtolength{\textwidth}{2cm}
\addtolength{\topmargin}{-2cm}
\addtolength{\textheight}{3.5cm}
\newcommand{\HRule}{\rule{\linewidth}{0.5mm}}
\makeindex

\usepackage{longtable}
\usepackage[pdftex]{graphicx}
\usepackage{makeidx}
\usepackage{hyperref}
\hypersetup{
    colorlinks=true,
    linkcolor=blue,
    filecolor=magenta,      
    urlcolor=cyan,
}


% define the title
\author{Men-at-Work}
\title{ OnlyRugby User Manuel: Demo 2}
\begin{document}
\setlength{\parskip}{6pt}

% generates the title
\begin{titlepage}

\begin{center}
% Upper part of the page       
\includegraphics[width=1\textwidth]{./up-logo.jpg}\\[0.4cm]    
\textsc{\LARGE Department of Computer Science}\\[1.5cm]
\textsc{\Large COS 301 - Software Engineering}\\[0.5cm]
% Title
\HRule \\[0.4cm]
{ \huge \bfseries OnlyRugby\\User Manuel: Demo 2}\\[0.4cm]
\HRule \\[0.4cm]
% Author and supervisor
\begin{minipage}{0.4\textwidth}
\begin{flushleft} \large
\emph{Authors:}
\end{flushleft}
\end{minipage}
\begin{minipage}{0.4\textwidth}
\begin{flushright} \large
\emph{Student number:}
\end{flushright}
\end{minipage}

\begin{minipage}{0.4\textwidth}
\begin{flushleft} \large
Herman {Keuris}
\end{flushleft}
\end{minipage}
\begin{minipage}{0.4\textwidth}
\begin{flushright} \large
\emph{}
u13037618
\end{flushright}
\end{minipage}

\begin{minipage}{0.4\textwidth}
\begin{flushleft} \large
Johan {van Rooyen}
\end{flushleft}
\end{minipage}
\begin{minipage}{0.4\textwidth}
\begin{flushright} \large
\emph{}
u11205131
\end{flushright}
\end{minipage}

\begin{minipage}{0.4\textwidth}
\begin{flushleft} \large
Estian {Rosslee}
\end{flushleft}
\end{minipage}
\begin{minipage}{0.4\textwidth}
\begin{flushright} \large
\emph{}
u12223426
\end{flushright}
\end{minipage}

\begin{minipage}{0.4\textwidth}
\begin{flushleft} \large
Ivan {Henning}
\end{flushleft}
\end{minipage}
\begin{minipage}{0.4\textwidth}
\begin{flushright} \large
\emph{}
u13008219
\end{flushright}
\end{minipage}

\begin{minipage}{0.4\textwidth}
\begin{flushleft} \large
Muller {Potgieter}
\end{flushleft}
\end{minipage}
\begin{minipage}{0.4\textwidth}
\begin{flushright} \large
\emph{}
u12003672
\end{flushright}
\end{minipage}

\vfill
% Bottom of the page
{\large \today}
\end{center}
\end{titlepage}
\footnotesize
%\input{declaration_of_originality.tex}
\normalsize


\pagenumbering{roman}
\tableofcontents
\newpage
\pagenumbering{arabic}

\newpage
\section{User Manuel}
\subsection{System Overview} 
The OnlyRugby App will serve school rugby teams and rugby clubs who have registered their teams with the main OnlyRugby website (implemented by our clients) by allowing them to log basic match time information. This means that when the user's team is playing a match the user can log such information as which team and players have scored, any injuries or penalties during the game and which team has won the line-out/scrum etc. This information is then stored on the main website's database allowing user's to compare statistics through the site.
\subsection{System Configuration}
\begin {itemize}
	\item The user will run the app itself from a phone using Android 4.0 or higher.
	\item All communication between the mobile app and the server will occur through HTTP and will consist of JSON objects which will be processed by a RESTful API on the server side (which will then format the information to be stored on the main OnlyRugby website's database) and the main app on the user's mobile device respectivley.
	\item The mobile app will also manage a small temporary SQLite database which will store any information waiting to be sent to the server when a connection is again established.
	\item The server will also use SMTP to send emails to an email account specified by the user in certain situations (e.g. when the user first registers an account with OnlyRugby). 
\end{itemize}
\begin{center}
  	 \includegraphics[width=0.5\textwidth] {./Sysyem Configuration.jpg}\\[0.4cm]
\end{center}
\subsection{Installation}
\begin {itemize}
	\item The user will have to choose to download the app from the Google App store and the rest will be done automatically by the app store itself. The user will be notified when the app finished downloading and is ready to run.
\end{itemize}

\subsection{Getting started}
	
\subsection{Using the system}
The functionality of the app will be spread between the following use cases:
	\subsection{Login/Logout}
		\begin {itemize}
			\item The first screen after running the App the user will be presented with two edit fields one for his username and one for his password, the user enters his details to login and continue using the app.
		\end{itemize}
	\subsection{Load Info}

	\subsection{Game Time}
	This use case is used to log the start and end time of each half of a match, as well as any intervals during which time was lost (the game was paused) and a reason for this time loss (injury, substitution, referee consultation, replacement of damaged player clothing). 
	\begin{itemize}Usage:
		\item The game time is displayed at the top of the UI.
		\item One starts the clock by tapping on it once. Preferably at kickoff. In the future the user will be presented with a confirm option in the form of a pop up which will not affect the running clock if accepted and reset the clock if denied.
		\item Once the clock is running a single tap on the clock will register a pause and activate a pop up asking the user for a reason for the pause, along with the legal options. The user will also have the option return should they accidentally activate this function. After a legal reason has been given for the pause the user is presented with a blank page and a button to resume the game.
		\item Upon reaching 40 minutes the game clock:
		\begin{itemize}
		\item registers the end of the current half,
		\item resets for the next half if there is one, and
		\item resets for overtime if conditions for it are met.
		\end{itemize}
	\end{itemize}

	\subsection{Scoring}
		During game time (i.e. while the OnlyRugby app is being used to log the information of a match currently being played ) the Scoring function will be used each time a player from either team scores points for their team (be it through a try, penalty kick, dropkick or a conversion kick following a try). The Scoring function works as follows:
		\begin{enumerate}
			\item Once the user has pressed the button to log score information they are redirected to a page which asks them to indicate which team scored by pressing either of two buttons which each have the names of the respective teams on them.
			\item The user is then redirected to a page which asks the user to specify what type of score it should log. This page displays buttons titled 'Try', 'Penalty Kick' and  'Dropkick' (note that 'Conversion Kick' is not an option as it can only be attempted after a succesful try). 
			\item The app then redirects the user to another page which has 15 buttons representing the 15 on-field players of the previously selected team. Each button has the player's position/jersey number on it as well as the player's name (if that information is provided on the database). This page also has an 'Unknown' button incase the user, for some reason, is not sure who scored on this team. The user is asked to select one of these players/buttons as the player who scored.
			\item In most cases the user is then done entering the score information but in the case of a 'Try' more steps must be completed:
				\begin{itemize}
					\item The user is first asked by a small pop-up message if the try was assisted by any other players with the options of clicking 'Yes' or 'No'. In the case of 'Yes' the user is then redirected to the same 'player selection page' as previously described were the user can then select a player which assisted the first player to score the try.
					\item The player is then asked by another small pop-up message if the subsequent conversion kick was successful with the options of clicking 'Yes' or 'No'. Selecting 'Yes' once again redirects the user to the 'player selection page' to select a player which successfully scored the conversion kick.
				\end{itemize}
		\end{enumerate}
	On all of the above described pages there are also buttons labeled 'Cancel' (to cancel logging the score) and 'Back' which redirects the user to the previous page and ignores the information that was logged on the page it is now leaving.

	\subsection{Substitutions}
	
	\subsection{Discipline}

	\subsection{Line-outs}
	Currently not implemented. Usage likely to be similar to the following description:
	\begin{itemize}
		\item Line out button present on landing page after login.
		\item Button can only be activated during game time.
		\item Activating the button registers a line out and redirects the interface to the line outs page. Here the user will be presented with options to specify:
			\begin{enumerate}
				\item The team to which the line out was given.
				\item The player to perform the throw in, by player number or by name. (Defaults to "unknown player").
				\item The team that won the line out.
			\end{enumerate}
		\item An option to return to the landing page will also be present
		\item Activating the return button will present a pop up prompting the user to save data as default, what the user entered, or to cancel the line out.
	\end{itemize}

	\subsection{Scrums}

	\subsection{Tackles}

	\subsection{Possession}

	\subsection{Turnovers}

	\subsection{Clean break}

	\subsection{Offloads}
	Currently not implemented. Usage likely to be similar to the following description:
	\begin{itemize}
		\item Offload button present on landing page after login.
		\item Button can only be activated during game time.
		\item Activating the button registers an offload and presents a pop up prompt asking for:
			\begin{enumerate}
				\item The team to which the offload was given. (2 buttons)
				\item The player that performed the offload, by player number or by name. (Defaults to "unknown player").
			\end{enumerate}
		\item An option to cancel the offload will also be present along with return option
		\item Activating the return button will prompt the user to save data as default, what the user entered, or to cancel the registering of an offload.

	\subsection{Ruck}

	\subsection{Maul}

\subsection{Troubleshooting}
		
\end{document}
